L'acide propanoïque $\ce{C2H5-COOH}$ est un acide gras, utilisé dans la synthèse
de certains produits organiques et pharmaceutiques, de parfums et dans la
médecine vétérinaire.

Le but de cette partie est l'étude de la réaction de l'acide propanoïque avec
l'eau et la détermination de la constante d'acidité du couple
$\ce{C2H5-COOH_{(aq)}}/\ce{C2H5-COO^{-}_{(aq)}}$.

\begin{enumerate}
    \item On considère, à $\qty{25}{\degreeCelsius}$, une solution aqueuse
          $(\mathrm{S})$ d'acide propanoïque de concentration molaire
          $C_{A}=\qty{2.0e-3}{\mol\per\L}$ et de volume $V_{A}=\qty{1.0}{\L}$.
          La mesure de la conductivité $\sigma$ de la solution $(\mathrm{S})$ a
          donné la valeur $\sigma=\qty{6.2e-3}{\S\per\m}$.

          \begin{donnees}
              \begin{itemize}
                  \item L'expression de la conductivité $\sigma$ de la solution
                        $(\mathrm{S})$~:
                        $\sigma=\lambda_{1}\bigl[\ce{H3O^{+}}\bigr]+
                        \lambda_{2}\bigl[\ce{C2H5-COO^{-}}\bigr]$, où
                        les concentrations sont exprimées en
                        ($\unit{\mol\per\cubic\m}$).
                  \item $\lambda_{1}=\lambda_{\ce{H3O^{+}}}=
                        \qty{35.0e-3}{\S\square\m\per\mol}$~;
                        $\lambda_{2}=\lambda_{\ce{C2H5COO^-}}=
                        \qty{3.58e-3}{\S\square\m\per\mol}$
              \end{itemize}
          \end{donnees}
          \begin{enumerate}
              \item Écrire l'équation chimique modélisant la réaction de l'acide
                    propanoïque avec l'eau.
              \item Dresser le tableau d'avancement de la réaction en utilisant
                    les grandeurs $C_{A}$, $V_{A}$, l'avancement $x$ et
                    l'avancement $x_{\eq}$ à l'état d'équilibre du système
                    chimique.
              \item Déterminer la valeur de l'avancement maximal $x_{\max}$.
              \item Vérifier que la valeur de l'avancement à l'état d'équilibre
                    est $x_{\eq}=\qty{1.6e-4}{\mol}$.
              \item Calculer la valeur du taux d'avancement final $\tau$.
                    Déduire.
              \item Vérifier que la valeur de la constante d'acidité du couple
                    $\ce{C2H5-COOH_{(aq)}}/\ce{C2H5-COO^{-}_{(aq)}}$ est
                    $\mathrm{K_{A}}\simeq\num{1.39e-5}$.
          \end{enumerate}

    \item On considère une solution aqueuse $(\mathrm{S^{\prime}})$ d'acide
          propanoïque de concentration molaire
          $C_{A}^{\prime}=\qty{2e-4}{\mol\per\L}$ et de $\mathrm{pH}=4,3$. On
          note $\tau^{\prime}$ le taux d'avancement final de la réaction de
          l'acide propanoïque avec l'eau dans ce cas.
          \begin{enumerate}
              \item Déterminer la valeur de $\tau^{\prime}$.
              \item Comparer les valeurs de $\tau$ et $\tau^{\prime}$. Déduire.
          \end{enumerate}
\end{enumerate}


